(1) 对绝热、无轴功、稳态流动控制体，单位质量稳流第一定律为 \(h_1+u_1^2/2+gz_1=h_2+u_2^2/2+gz_2\)。忽略流速与高度变化后化为 \(h_1=h_2\)，并非 \(u_1=u_2\)；压强降低时流动功已包含在焓 \(h=u+pv\) 中。实际节流通常有熵产生，故不能另外假定等熵。
(2) 由 Maxwell 关系 \((\partial s/\partial p)_T=-(\partial v/\partial T)_p\)，有 \(dh=C_p\,dT+[v-T(\partial v/\partial T)_p]\,dp\)。令 \(dh=0\)，立刻得到 \(\mu_{JT}=(\partial T/\partial p)_h=[T(\partial v/\partial T)_p-v]/C_p\)。系数大于零意味着压强降低 \(dp<0\) 时温度降低，注意符号不能倒置。
(3) 理想气体 \(v=RT/p\)，所以 \(T(\partial v/\partial T)_p-v=0\)，节流不改变温度。在范德瓦耳斯模型低压区，以摩尔体积展开 \(v\simeq RT/p+b-a/(RT)\)，从而 \((\partial v/\partial T)_p\simeq R/p+a/(RT^2)\)。代入得 \(\mu_{JT}\simeq[2a/(RT)-b]/C_p\)，省略的是较高阶的压强修正。
(4) 低压近似下温度变化的符号在 \(2a/(RT)-b=0\) 时翻转，即 \(T_{\rm inv}\simeq2a/(Rb)\)。低于此温度节流降温，高于此温度节流升温；这是 \(dp<0\) 情形。真实高压下 \(v\) 对压强的非线性不可忽略，可能存在完整逆转曲线及相变，因此不能把这一单一数值当作所有压强处的精确液化阈值。
